% Additional arm structure; not the general one-descent theorem.
\section{The opposite reflection for increasing-arm deficits}
\begin{theorem}
For a nonempty increasing arm of weight $a$, every height layer
of $E_X=1-z+(2z-1)F_X$ has endpoint sum at least $a+3$.
Equivalently, every layer of its reverse $e_X=z^{a+1}E_X(1/z)$
has endpoint sum at most $a-1$.
\end{theorem}
\begin{proof}
For a single part of size $a$, or any interval arm polynomial,
$E_X=z^2[a-1]+2z^{a+1}$, whose layers have the asserted centers.
Induct under appending a current maximum $t$. Write $c$ for
the old arm weight and $E$ for its deficit. The exact recurrence
is $E_{Zt}=W+z^{t+1}$, where $W=S_tE$.
The old layer bound makes every layer-product endpoint sum of
$W$ at least $c+t+4$. Hence $W$ satisfies reflection about
$(c+t+3)/2$ with one unit of center margin.

If $c<t$, the added singleton lies on the right of that center
(or at the center), so it preserves the reflected inequalities.
Suppose $c\ge t$. Only the reflected pair at $t+1,c+2$ can
lose one. We prove $W_{c+2}-W_{t+1}\ge1$.
The old reflection gives, after pairing indices beyond its midpoint,
\[
 \sum_{i=1}^{t}E_i\le\sum_{i=c-t+3}^{c+2}E_i.
\]
Since $E_{c+2}=0$, this implies
\[
 W_{c+2}-W_{t+1}
 =\sum_{i=c-t+2}^{c+1}E_i-\sum_{i=1}^{t}E_i
 \ge E_{c-t+2}.
\]
If $c-t+2$ belongs to the positive support $[\kappa,c+1]$,
this lower bound is at least one. Otherwise $c-t+2<\kappa$,
so the first sum contains the entire deficit. The second sum
omits the final coefficient $E_{c+1}=2$, because $t\le c$.
The difference is then at least two.
Thus the extra singleton is paid for in every case.
The increasing-arm deficit is already known to be unimodal,
so the reflected inequalities imply the stated height-layer bound.
\end{proof}

Together with the reverse lower bounds proved separately, a
normalized nonrectangular arm $1^\alpha U$, with
$s=\operatorname{len}U$, has deficit-layer endpoint sums in
$[s,a-1]$ after reversal. The lower endpoint improves to $s+1$
unless the arm is $1,2^s$.

\paragraph{Sharpness of the uniform forward constant.}
It is false that every nonrectangular increasing arm has deficit
layer endpoint sums at least $a+4$. For
$X=1^9(6,18,19,38,91)$ one has $a=181$ and the minimum sum is
$184=a+3$. This counterexample does not affect the theorem.

Nor should the two-sided layer bounds be replaced by log-concavity
of the reverse deficit on its increasing side. For $X=(1,2,10)$
the first four reverse-deficit coefficients are $2,5,5,6$;
the log-concavity inequality at the third coefficient fails.
